\zag{Źródła prawa karnego procesowego}[{\kolor A}]\str{52-55}

Najważniejszymi źródłami prawa karnego procesowego są \acr{kpk}, którego \textit{ratio} jest regulowanie tegoż procesu, oraz \acr{krp}, przyznająca jednostkom gwarancje i prawa mające zastosowanie w procesie karnym i regulująca zrębowo podstawowe instytucje prawa karnego procesowego.\medskip

Ważnymi źródłami tego prawa są także akty prawa międzynarodowego, wśród których wymienić należy:\begin{itemize}
    \item \acr{ekpcz};
    \item \acr{mppiop};
    \item \acr{kppue};
    \item dyrektywy:\begin{enumerate}
        \item informacyjną,
        \item o dostępie do adwokata,
        \item niewinnościową,
        \item o pomocy prawnej.
    \end{enumerate}
\end{itemize}

Regulacje \acr{kpk} uzupełniają także liczne polskie ustawy, które można podzielić na trzy grupy (wypisane skrótami, żeby niepotrzebnie nie powielać treści wykazu skrótów):\begin{enumerate}
    \item regulujące pozakodeksowe tryby szczególne postępowania karnego:\\ \acr{kks}, \acr{odp.podm.zbior};
    \item regulujące ustrój i funkcjonowanie organów i innych uczestników postępowania karnego:\\ \acr{uosn}, \acr{pusp}, \acr{pp}, \acr{uop}, \acr{padw}, \acr{uorp}, \acr{uoabw}, \acr{uocba}, \acr{uosg}, \acr{uorpo};
    \item regulujące inne kwestie szczególne:\\ \acr{uoipn}, \acr{rep.niep}, \acr{uopn}, \acr{uo60}.
    
\end{enumerate} 